\section{An exact one-bit radial repair}
\label{supp:one-bit}

At displacement $h=1/2$, take the midpoint $c=(1/2,1/4)$ and write $v=x-c$. Retain the direction $u=v/\|v\|$ and the binary radial code
\begin{equation}
 q(x)=\mathbf 1\{\|v\|>4/5\}.
 \label{supp:eq:one-bit-code}
\end{equation}
The following decoder gives an explicit decision-preserving representation:
\begin{equation}
 \widehat y(u,q)=
 \begin{cases}
 0,&q=0\text{ and }9u_1+10u_2>0,\\
 0,&q=1\text{ and }u_2-u_1>0,\\
 1,&\text{otherwise}.
 \end{cases}
 \label{supp:eq:one-bit-decoder}
\end{equation}

\paragraph{Proposition S1 (minimal radial repair at the midpoint).}
For every $x\in S_0\cup S_1$ at $h=1/2$, the decoder in~\eqref{supp:eq:one-bit-decoder} returns its class label. Thus $(T_c(X),q(X))$ has zero Bayes risk for either study sampling law, and more generally for arbitrary densities positive almost everywhere on these disjoint supports and arbitrary positive class priors. Exactly one additional radial bit is sufficient and minimal among finite-valued radial codes appended to exact direction.

\emph{Proof.}
It suffices to establish two separating inequalities for $S_0$; $S_1=(1,1/2)-S_0$ is its reflection about $c$, which preserves the radius and reverses every linear form in $v$.

First consider the near region $\|v\|\leq4/5$. Suppose also that $9v_1+10v_2\leq0$, or $9x_1+10x_2\leq7$. Intersect these inequalities with $x_2\geq0$. The resulting compact convex region has a circular boundary arc and two straight boundary segments. On a straight segment the convex function $\|x\|^2$ attains its maximum at an endpoint. Parametrize the circular arc by
\[
 x=c+\tfrac45(\cos\phi,\sin\phi),\qquad
 \operatorname{atan2}(9,-10)\leq\phi\leq\pi+\arcsin(5/16).
\]
Its squared norm is
\[
 \|x\|^2=\frac{381}{400}
        +\frac85\left(\frac12\cos\phi+\frac14\sin\phi\right).
\]
This expression decreases along the displayed arc: its two stationary angles are $\arctan(1/2)$ and $\pi+\arctan(1/2)$, both outside the arc. The largest circular-arc value is therefore
\[
 \frac{381}{400}-\frac{22}{5\sqrt{181}}<\frac{16}{25}.
\]
The remaining relevant vertex is $(7/9,0)$, whose squared norm is $49/81<16/25$. The circular endpoint on $x_2=0$ has a smaller norm than the upper arc endpoint. The entire forbidden region thus lies inside the class-0 inner disk. Since $S_0$ requires $\|x\|>4/5$, every near point in $S_0$ satisfies $9v_1+10v_2>0$.

For the far region, suppose instead that $v_2-v_1\leq0$, or $x_1-x_2\geq1/4$. The compact convex region
\[
 \{x:x_2\geq0,\ x_1-x_2\geq1/4,\ \|x\|\leq6/5\}
\]
consists of an outer-circle arc and two straight segments. Its noncircular vertex $(1/4,0)$ has squared distance $1/8$ from $c$. On the circular arc write $x=(6/5)(\cos\theta,\sin\theta)$; its angles satisfy $0\leq\theta<\pi/4$. Throughout that interval $\cos\theta+(1/2)\sin\theta\geq1$, so
\[
 \|x-c\|^2=\frac{701}{400}-x_1-\frac12x_2
 \leq\frac{221}{400}<\frac{16}{25}.
\]
Convexity bounds each straight segment by its endpoint values. Thus every point of the forbidden region is near, and every far point in $S_0$ satisfies $v_2-v_1>0$. Dividing these strict separating inequalities by the positive radius proves the decoder. Reflection proves the class-1 assertions.

The direction-only representation has strictly positive Bayes risk by the all-origin result. A radial code with no additional bit has only one value and retains precisely that representation. The binary code above achieves zero risk, establishing the stated minimality.\hfill$\square$
